\section{Routine facts}\label{app:routine}

This appendix proves the routine facts used in Sections~\ref{sec:setting} to~\ref{sec:local}.

\subsection{Regular variation}

\begin{lemma}\label{lem:rv-facts}
Let $T$ be regularly varying with index $-\alpha\in(-1,0)$, nonincreasing and positive.
\begin{enumerate}[(i)]
\item (R4) of Section~\ref{sec:rv} holds.
\item If $a$ is increasing with $tT(a(t))\to1$, then $\sup_{m\ge1}mT(a_m)<\infty$.
\item With $m_N(t)=\lfloor t/T(N)\rfloor$, $a_{m_N(t)}/N\to t^{1/\alpha}$ uniformly for $t$ in compact
subsets of $(0,\infty)$.
\end{enumerate}
\end{lemma}

\begin{proof}
(i) The function $\phi(w)=wT(w)$ is regularly varying with index $1-\alpha>0$, and it is
positive, so $\inf_{w\ge1}\phi(w)>0$ because $\phi(w)\to\infty$. By Potter's
theorem~\cite[Theorem~1.5.6]{bgt1989} there is
$x_0$ with $\phi(y)/\phi(x)\le2(y/x)^{(1-\alpha)/2}$ for $x_0\le y\le x$. By (R2),
$2^lT(2^{l-1})\le C\phi(2^l)$. The terms with $2^l\ge x_0$ add up to at most
$C\phi(x)\sum_{l\ge0}2^{-l(1-\alpha)/2}$, and the finitely many terms with $2^l<x_0$ add up to a
constant, which is at most $C\phi(x)$. For the second sum let $x\ge x_0$. The first term
with $2^l>x$ has $2^{l-1}\in(x/2,x]$, so it is at most $T(x/2)\le CT(x)$ by (R2). The later
terms have $2^{l-1}\ge x$, and Potter's bounds give $T(y)/T(x)\le2(y/x)^{-\alpha/2}$ for
$y\ge x\ge x_0$. So the sum is at most $CT(x)\sum_{l\ge0}2^{-l\alpha/2}$. For $1\le x<x_0$ the sum is at most
$\sum_{l\ge1}T(2^{l-1})<\infty$, which is at most $CT(x_0)\le CT(x)$.

(ii) The sequence $mT(a_m)$ converges to $1$ and each term is finite.

(iii) Fix $t>0$ and write $m=m_N(t)$. Then $mT(N)\to t$ and $mT(a_m)\to1$, so
$T(a_m)/T(N)\to1/t$. Take a subsequence along which $a_m/N\to\lambda\in[0,\infty]$. If
$\lambda\in(0,\infty)$, (R1) gives $T(a_m)/T(N)\to\lambda^{-\alpha}$, so $\lambda=t^{1/\alpha}$. If
$\lambda=0$ or $\lambda=\infty$, Potter's bounds give $T(a_m)/T(N)\to\infty$ or $\to0$, which is
impossible. So $a_m/N\to t^{1/\alpha}$. The functions $t\mapsto a_{m_N(t)}/N$ are nondecreasing and
the limit is continuous and increasing, so the convergence is uniform on compact sets.
\end{proof}

\subsection{The Poisson integral}

\begin{lemma}\label{lem:poisson}
Let $\Lambda\ge0$ be integrable on $(0,1)$ and continuous at $x\in(0,1)$, and let
$G(z)=\int_0^1\Lambda(v)\,dv/(z-v)$ for $z\notin[0,1]$. Then
$-\frac1\pi\mathrm{Im}\,G(x+i\varepsilon)\to\Lambda(x)$ as $\varepsilon\downarrow0$.
\end{lemma}

\begin{proof}
Since $\mathrm{Im}\,(x+i\varepsilon-v)^{-1}=-\varepsilon/((x-v)^2+\varepsilon^2)$,
\[
 -\frac1\pi\,\mathrm{Im}\,G(x+i\varepsilon)=\int_0^1\Lambda(v)\,P_\varepsilon(x-v)\,dv,\qquad
 P_\varepsilon(u)=\frac{\varepsilon}{\pi(u^2+\varepsilon^2)} .
\]
The kernel $P_\varepsilon$ is positive with total mass $1$ on $\R$, and its mass outside
$[-\eta,\eta]$ and its supremum there tend to $0$ for each $\eta>0$. Split the integral at
$|x-v|\le\eta$. The near part differs from $\Lambda(x)\int_{|u|\le\eta}P_\varepsilon$ by at most
$\sup_{|v-x|\le\eta}|\Lambda(v)-\Lambda(x)|$, and the far part is at most
$\sup_{|u|>\eta}P_\varepsilon(u)\int_0^1\Lambda$. Let $\varepsilon\to0$ and then $\eta\to0$.
\end{proof}

\subsection{The normalization of the minus sums}

Under (RV$_\alpha$), $mT_\minus(a_mz)\to bz^{-\alpha}$ for every $z>0$, by (R1) and
$mT_\plus(a_m)\to1$. So the sum of $m$ minus runs, divided by $a_m$, converges in law to the
positive stable variable whose L\'evy measure has tail $bz^{-\alpha}$, which is $b^{1/\alpha}Z$.
Gnedenko's local limit theorem~\cite[Theorem~4.2.1]{ibragimovlinnik1971} holds for every
norming sequence under which the sums converge, so it applies with the norming $a_m$ and the
density $g_b$. This is the form used in Lemma~\ref{lem:cd}(i).
